% ============================================================
%  PROBLEMAS RUTINARIOS
%  Michael → Actividad 1
%  David   → Actividad 2
%  Fancho  → Actividad 3
% ============================================================

\section{Problemas Rutinarios:}

Se espera que las siguientes actividades le permitan comprender, repasar y utilizar los
métodos analíticos: factor integrante, variables separables y ecuaciones exactas (entre
otros). Además, que el desarrollo de estas actividades genere en usted sentimientos de
confianza al momento de resolver algunas ecuaciones diferenciales ordinarias de primer
orden.

% ──────────────────────────────────────────────
%  ACTIVIDAD 1  →  Michael
% ──────────────────────────────────────────────

\actividad{1} Realizar los ejercicios del 1 al 16 de la página 33 de [2], entregar
4, 8, 10, 12, 13, 14 y 16.

% ── Ejercicio 4 ─────────────────────────────
\subsection*{Ejercicio 4 \quad\normalfont\small(pág. 33, [2])}

\begin{enunciado}{Enunciado}
    % Copiar/transcribir el enunciado del ejercicio 4 de la página 33 de Boyce-DiPrima.
    \textit{[Transcribir el enunciado del ejercicio 4 aquí.]}
\end{enunciado}

\begin{respuesta}
    % Michael — desarrollo y solución
    \vspace{7cm}
\end{respuesta}

% ── Ejercicio 8 ─────────────────────────────
\subsection*{Ejercicio 8 \quad\normalfont\small(pág. 33, [2])}

\begin{enunciado}{Enunciado}
    \textit{[Transcribir el enunciado del ejercicio 8 aquí.]}
\end{enunciado}

\begin{respuesta}
    % Michael — desarrollo y solución
    \vspace{7cm}
\end{respuesta}

% ── Ejercicio 10 ────────────────────────────
\subsection*{Ejercicio 10 \quad\normalfont\small(pág. 33, [2])}

\begin{enunciado}{Enunciado}
    \textit{[Transcribir el enunciado del ejercicio 10 aquí.]}
\end{enunciado}

\begin{respuesta}
    % Michael — desarrollo y solución
    \vspace{7cm}
\end{respuesta}

% ── Ejercicio 12 ────────────────────────────
\subsection*{Ejercicio 12 \quad\normalfont\small(pág. 33, [2])}

\begin{enunciado}{Enunciado}
    \textit{[Transcribir el enunciado del ejercicio 12 aquí.]}
\end{enunciado}

\begin{respuesta}
    % Michael — desarrollo y solución
    \vspace{7cm}
\end{respuesta}

% ── Ejercicio 13 ────────────────────────────
\subsection*{Ejercicio 13 \quad\normalfont\small(pág. 33, [2])}

\begin{enunciado}{Enunciado}
    \textit{[Transcribir el enunciado del ejercicio 13 aquí.]}
\end{enunciado}

\begin{respuesta}
    % Michael — desarrollo y solución
    \vspace{7cm}
\end{respuesta}

% ── Ejercicio 14 ────────────────────────────
\subsection*{Ejercicio 14 \quad\normalfont\small(pág. 33, [2])}

\begin{enunciado}{Enunciado}
    \textit{[Transcribir el enunciado del ejercicio 14 aquí.]}
\end{enunciado}

\begin{respuesta}
    % Michael — desarrollo y solución
    \vspace{7cm}
\end{respuesta}

% ── Ejercicio 16 ────────────────────────────
\subsection*{Ejercicio 16 \quad\normalfont\small(pág. 33, [2])}

\begin{enunciado}{Enunciado}
    \textit{[Transcribir el enunciado del ejercicio 16 aquí.]}
\end{enunciado}

\begin{respuesta}
    % Michael — desarrollo y solución
    \vspace{7cm}
\end{respuesta}

\newpage

% ──────────────────────────────────────────────
%  ACTIVIDAD 2  →  David
% ──────────────────────────────────────────────

\actividad{2} Realizar los ejercicios del 1 al 16 de la página 40 de [2], entregar
8, 10, 12, 14 y 16.

% ── Ejercicio 8 ─────────────────────────────
\subsection*{Ejercicio 8 \quad\normalfont\small(pág. 40, [2])}

\begin{enunciado}{Enunciado}
    \textit{[Transcribir el enunciado del ejercicio 8 aquí.]}
\end{enunciado}

\begin{respuesta}
    % David — desarrollo y solución
    \vspace{7cm}
\end{respuesta}

% ── Ejercicio 10 ────────────────────────────
\subsection*{Ejercicio 10 \quad\normalfont\small(pág. 40, [2])}

\begin{enunciado}{Enunciado}
    \textit{[Transcribir el enunciado del ejercicio 10 aquí.]}
\end{enunciado}

\begin{respuesta}
    % David — desarrollo y solución
    \vspace{7cm}
\end{respuesta}

% ── Ejercicio 12 ────────────────────────────
\subsection*{Ejercicio 12 \quad\normalfont\small(pág. 40, [2])}

\begin{enunciado}{Enunciado}
    \textit{[Transcribir el enunciado del ejercicio 12 aquí.]}
\end{enunciado}

\begin{respuesta}
    % David — desarrollo y solución
    \vspace{7cm}
\end{respuesta}

% ── Ejercicio 14 ────────────────────────────
\subsection*{Ejercicio 14 \quad\normalfont\small(pág. 40, [2])}

\begin{enunciado}{Enunciado}
    \textit{[Transcribir el enunciado del ejercicio 14 aquí.]}
\end{enunciado}

\begin{respuesta}
    % David — desarrollo y solución
    \vspace{7cm}
\end{respuesta}

% ── Ejercicio 16 ────────────────────────────
\subsection*{Ejercicio 16 \quad\normalfont\small(pág. 40, [2])}

\begin{enunciado}{Enunciado}
    \textit{[Transcribir el enunciado del ejercicio 16 aquí.]}
\end{enunciado}

\begin{respuesta}
    % David — desarrollo y solución
    \vspace{7cm}
\end{respuesta}

\newpage

% ──────────────────────────────────────────────
%  ACTIVIDAD 3  →  Fancho
% ──────────────────────────────────────────────

\actividad{3} Realizar los ejercicios del 1 al 22 de la página 77 de [2], entregar
10, 12, 14, 16, 18 y 22.

% ── Ejercicio 10 ────────────────────────────
\subsection*{Ejercicio 10 \quad\normalfont\small(pág. 77, [2])}

\begin{enunciado}{Enunciado}
    Resolver el siguiente problema con valor inicial y determinar 
    al menos aproximadamente dónde la solución es válida:
    \[
        (9x^2 + y - 1) - (4y - x)y' = 0, \quad y(1) = 0
    \]
\end{enunciado}

\begin{respuesta}

\textbf{Paso 1: Identificar el tipo de ecuación.}

La ecuación tiene la forma $M(x,y) + N(x,y)y' = 0$, donde:
\[
    M = 9x^2 + y - 1 \qquad N = -(4y - x) = x - 4y
\]
Para verificar si es exacta, calculamos las derivadas parciales cruzadas:
\[
    \frac{\partial M}{\partial y} = 1 \qquad \frac{\partial N}{\partial x} = 1
\]
Como $\dfrac{\partial M}{\partial y} = \dfrac{\partial N}{\partial x}$, 
la ecuación \textbf{es exacta}.

\medskip
\textbf{Paso 2: Encontrar la función $F(x,y)$ tal que $F_x = M$ y $F_y = N$.}

Integramos $M$ respecto a $x$:
\[
    F = \int (9x^2 + y - 1)\, dx = 3x^3 + xy - x + h(y)
\]
donde $h(y)$ es una función desconocida que solo depende de $y$, 
ya que al derivar respecto a $x$ cualquier función pura de $y$ desaparece.

\medskip
\textbf{Paso 3: Determinar $h(y)$ usando la segunda condición $F_y = N$.}

Derivamos $F$ respecto a $y$ e igualamos a $N$:
\[
    \frac{\partial F}{\partial y} = x + h'(y) = x - 4y
\]
Despejando:
\[
    h'(y) = -4y \implies h(y) = -2y^2
\]

\medskip
\textbf{Paso 4: Escribir la solución general.}

Juntando todo:
\[
    F(x,y) = 3x^3 + xy - x - 2y^2 = C
\]

\medskip
\textbf{Paso 5: Aplicar la condición inicial $y(1) = 0$.}

Sustituimos $x = 1$, $y = 0$:
\[
    3(1)^3 + (1)(0) - (1) - 2(0)^2 = C \implies C = 2
\]
La solución del problema con valor inicial es:
\[
    \boxed{3x^3 + xy - x - 2y^2 = 2}
\]

\medskip
\textbf{Paso 6: Validez de la solución.}

Como $M$ y $N$ son polinomios, son continuas en todo $\mathbb{R}^2$. 
Por el Teorema de Existencia y Unicidad, la solución es válida 
en algún intervalo que contiene $x = 1$. Dado que no hay 
singularidades, la solución es válida para todo $x \in \mathbb{R}$, 
aunque la relación implícita puede no definir $y$ como función 
explícita en todo ese dominio.

\end{respuesta}

% ── Ejercicio 12 ────────────────────────────
\subsection*{Ejercicio 12 \quad\normalfont\small(pág. 77, [2])}
\begin{enunciado}{Enunciado}
    Encontrar el valor de $b$ para el cual la siguiente ecuación es exacta,
    y luego resolverla:
    \[
        \left(ye^{2xy} + x\right) + bxe^{2xy}y' = 0
    \]
\end{enunciado}

\begin{respuesta}

\textbf{Paso 1: Identificar $M$ y $N$.}

La ecuación tiene la forma $M + Ny' = 0$, donde:
\[
    M = ye^{2xy} + x \qquad N = bxe^{2xy}
\]

\medskip
\textbf{Paso 2: Encontrar $b$ usando la condición de exactitud.}

Para que la ecuación sea exacta se debe cumplir 
$\dfrac{\partial M}{\partial y} = \dfrac{\partial N}{\partial x}$.

Calculamos cada parcial:
\[
    \frac{\partial M}{\partial y} = e^{2xy} + y \cdot 2x e^{2xy} 
    = e^{2xy}(1 + 2xy)
\]
\[
    \frac{\partial N}{\partial x} = be^{2xy} + bx \cdot 2ye^{2xy} 
    = be^{2xy}(1 + 2xy)
\]

Igualando:
\[
    e^{2xy}(1 + 2xy) = be^{2xy}(1 + 2xy) \implies \boxed{b = 1}
\]

\medskip
\textbf{Paso 3: Reescribir la ecuación con $b = 1$.}
\[
    \left(ye^{2xy} + x\right) + xe^{2xy}y' = 0
\]

\medskip
\textbf{Paso 4: Encontrar $F(x,y)$ tal que $F_x = M$ y $F_y = N$.}

Integramos $N = xe^{2xy}$ respecto a $y$:
\[
    F = \int xe^{2xy}\, dy = \frac{1}{2}e^{2xy} + g(x)
\]
donde $g(x)$ es una función desconocida que solo depende de $x$.

\medskip
\textbf{Paso 5: Determinar $g(x)$ usando la condición $F_x = M$.}

Derivamos $F$ respecto a $x$ e igualamos a $M$:
\[
    \frac{\partial F}{\partial x} = ye^{2xy} + g'(x) = ye^{2xy} + x
\]
Despejando:
\[
    g'(x) = x \implies g(x) = \frac{x^2}{2}
\]

\medskip
\textbf{Paso 6: Escribir la solución general.}

Juntando todo:
\[
    F(x,y) = \frac{1}{2}e^{2xy} + \frac{x^2}{2} = C
\]

Multiplicando por 2:
\[
    \boxed{e^{2xy} + x^2 = C}
\]

\end{respuesta}

% ── Ejercicio 14 ────────────────────────────
\subsection*{Ejercicio 14 \quad\normalfont\small(pág. 77, [2])}

\begin{enunciado}{Enunciado}
    Mostrar que cualquier ecuación separable de la forma
    \[
        M(x) + N(y)\,y' = 0
    \]
    es también una ecuación exacta.
\end{enunciado}

\begin{respuesta}

\textbf{Planteamiento.}

Una ecuación de la forma $M(x) + N(y)y' = 0$ es separable porque 
$M$ depende \textit{solo} de $x$ y $N$ depende \textit{solo} de $y$.

\medskip
\textbf{Verificación de la condición de exactitud.}

Para que sea exacta se debe cumplir:
\[
    \frac{\partial M}{\partial y} = \frac{\partial N}{\partial x}
\]

Como $M$ depende \textbf{solo de $x$}:
\[
    \frac{\partial M}{\partial y} = 0
\]

Como $N$ depende \textbf{solo de $y$}:
\[
    \frac{\partial N}{\partial x} = 0
\]

Por lo tanto:
\[
    \frac{\partial M}{\partial y} = 0 = \frac{\partial N}{\partial x} \checkmark
\]

La condición de exactitud se cumple para \textbf{cualquier} $M(x)$ 
y $N(y)$, independientemente de su forma. Luego, toda ecuación 
separable de esta forma es también exacta. $\blacksquare$

\medskip
\textbf{Nota.} La verificación no depende de los valores específicos 
de $M$ o $N$, sino únicamente de que cada función dependa de una 
sola variable. Eso garantiza que sus derivadas parciales respecto 
a la otra variable sean siempre cero.

\end{respuesta}
% ── Ejercicio 16 ────────────────────────────
\subsection*{Ejercicio 16 \quad\normalfont\small(pág. 77, [2])}
\begin{enunciado}{Enunciado}
    Mostrar que la siguiente ecuación no es exacta, pero que se 
    vuelve exacta al multiplicarla por el factor integrante dado.
    Luego resolver la ecuación:
    \[
        (x + 2)\sin(y) + (x\cos(y))\,y' = 0, 
        \qquad \mu(x,y) = xe^x
    \]
\end{enunciado}

\begin{respuesta}

\textbf{Paso 1: Verificar que la ecuación original NO es exacta.}

Identificamos $M = (x+2)\sin(y)$ y $N = x\cos(y)$:
\[
    \frac{\partial M}{\partial y} = (x+2)\cos(y) \qquad 
    \frac{\partial N}{\partial x} = \cos(y)
\]
Como $(x+2)\cos(y) \neq \cos(y)$, la ecuación \textbf{no es exacta}.

\medskip
\textbf{Paso 2: Multiplicar por el factor integrante $\mu = xe^x$.}

\[
    xe^x(x+2)\sin(y) + x^2e^x\cos(y)\,y' = 0
\]
Ahora: $M^* = xe^x(x+2)\sin(y)$ y $N^* = x^2e^x\cos(y)$.

\medskip
\textbf{Paso 3: Verificar exactitud de la ecuación modificada.}

\[
    \frac{\partial M^*}{\partial y} = xe^x(x+2)\cos(y)
\]
\[
    \frac{\partial N^*}{\partial x} = (2xe^x + x^2e^x)\cos(y) 
    = xe^x(x+2)\cos(y)
\]
Como $\dfrac{\partial M^*}{\partial y} = \dfrac{\partial N^*}{\partial x}$,
la ecuación modificada \textbf{es exacta}. $\checkmark$

\medskip
\textbf{Paso 4: Encontrar $F(x,y)$ tal que $F_x = M^*$ y $F_y = N^*$.}

Integramos $N^*$ respecto a $y$:
\[
    F = \int x^2e^x\cos(y)\,dy = x^2e^x\sin(y) + g(x)
\]

\medskip
\textbf{Paso 5: Determinar $g(x)$ usando $F_x = M^*$.}

\[
    \frac{\partial F}{\partial x} = (2xe^x + x^2e^x)\sin(y) + g'(x) 
    = xe^x(x+2)\sin(y) + g'(x)
\]
Igualando a $M^*$:
\[
    g'(x) = 0 \implies g(x) = 0
\]

\medskip
\textbf{Paso 6: Solución general.}

\[
    \boxed{x^2e^x\sin(y) = C}
\]

\medskip
\textbf{Verificación.} Derivando implícitamente respecto a $x$:
\[
    (2xe^x + x^2e^x)\sin(y) + x^2e^x\cos(y)\,y' = 0
\]
\[
    xe^x(x+2)\sin(y) + x^2e^x\cos(y)\,y' = 0
\]
Dividiendo entre $xe^x$:
\[
    (x+2)\sin(y) + x\cos(y)\,y' = 0 \checkmark
\]

\end{respuesta}

% ── Ejercicio 18 ────────────────────────────
\subsection*{Ejercicio 18 \quad\normalfont\small(pág. 77, [2])}
\begin{enunciado}{Enunciado}
    Encontrar un factor integrante y resolver la ecuación:
    \[
        (3x^2y + 2xy + y^3) + (x^2 + y^2)\,y' = 0
    \]
\end{enunciado}

\begin{respuesta}

\textbf{Paso 1: Verificar que la ecuación NO es exacta.}

Identificamos $M = 3x^2y + 2xy + y^3$ y $N = x^2 + y^2$:
\[
    \frac{\partial M}{\partial y} = 3x^2 + 2x + 3y^2 \qquad
    \frac{\partial N}{\partial x} = 2x
\]
Como $3x^2 + 2x + 3y^2 \neq 2x$, la ecuación \textbf{no es exacta}.

\medskip
\textbf{Paso 2: Encontrar el factor integrante.}

Calculamos el cociente $\dfrac{M_y - N_x}{N}$:
\[
    \frac{M_y - N_x}{N} 
    = \frac{(3x^2 + 2x + 3y^2) - 2x}{x^2 + y^2}
    = \frac{3(x^2 + y^2)}{x^2 + y^2} = 3
\]
Como el resultado depende \textbf{solo de $x$}, el factor integrante es:
\[
    \mu(x) = e^{\int 3\,dx} = e^{3x}
\]

\medskip
\textbf{Paso 3: Multiplicar por $\mu = e^{3x}$ y verificar exactitud.}

\[
    M^* = e^{3x}(3x^2y + 2xy + y^3) \qquad N^* = e^{3x}(x^2 + y^2)
\]
\[
    \frac{\partial M^*}{\partial y} = e^{3x}(3x^2 + 2x + 3y^2)
\]
\[
    \frac{\partial N^*}{\partial x} = 3e^{3x}(x^2+y^2) + e^{3x}(2x) 
    = e^{3x}(3x^2 + 3y^2 + 2x) \checkmark
\]
La ecuación modificada \textbf{es exacta}.

\medskip
\textbf{Paso 4: Encontrar $F(x,y)$ integrando $F_x = M^*$ respecto a $x$.}

Dado que $M^* = e^{3x}(3x^2y + 2xy + y^3)$, integramos respecto a $x$
agrupando por partes:
\[
    F = \int e^{3x}\bigl[(3x^2 + 2x)y + y^3\bigr]\,dx
\]

Para $\displaystyle\int(3x^2 + 2x)e^{3x}\,dx$, observamos que 
$(x^2 e^{3x})' = 2xe^{3x} + 3x^2e^{3x} = (3x^2+2x)e^{3x}$,
por lo que:
\[
    \int(3x^2+2x)e^{3x}\,dx = x^2e^{3x}
\]

Para $\displaystyle\int y^3 e^{3x}\,dx = \dfrac{y^3 e^{3x}}{3}$.

Entonces:
\[
    F = e^{3x}\!\left(x^2 y + \frac{y^3}{3}\right) + h(y)
\]

\medskip
\textbf{Paso 5: Determinar $h(y)$ usando $F_y = N^*$.}

\[
    \frac{\partial F}{\partial y} = e^{3x}\!\left(x^2 + y^2\right) 
    + h'(y) = N^* = e^{3x}(x^2 + y^2)
\]
\[
    h'(y) = 0 \implies h(y) = 0
\]

\medskip
\textbf{Paso 6: Solución general.}

\[
    \boxed{e^{3x}\!\left(x^2 y + \frac{y^3}{3}\right) = C}
\]

\medskip
\textbf{Verificación.} Derivando implícitamente respecto a $x$:
\[
    3e^{3x}\!\left(x^2y + \frac{y^3}{3}\right) 
    + e^{3x}\!\left(2xy + (x^2 + y^2)y'\right) = 0
\]
Dividiendo entre $e^{3x}$:
\[
    3x^2y + y^3 + 2xy + (x^2 + y^2)y' = 0
\]
\[
    (3x^2y + 2xy + y^3) + (x^2 + y^2)y' = 0 \checkmark
\]

\end{respuesta}
% ── Ejercicio 22 ────────────────────────────
\subsection*{Ejercicio 22 \quad\normalfont\small(pág. 77, [2])}
\begin{enunciado}{Enunciado}
    Resolver la siguiente ecuación diferencial usando el factor 
    integrante dado:
    \[
        (3xy + y^2) + (x^2 + xy)\,y' = 0, 
        \qquad \mu(x,y) = \frac{1}{xy(2x+y)}
    \]
\end{enunciado}

\begin{respuesta}

\textbf{Paso 1: Verificar que la ecuación original NO es exacta.}

Identificamos $M = 3xy + y^2$ y $N = x^2 + xy$:
\[
    \frac{\partial M}{\partial y} = 3x + 2y \qquad 
    \frac{\partial N}{\partial x} = 2x + y
\]
Como $3x + 2y \neq 2x + y$, la ecuación \textbf{no es exacta}.

\medskip
\textbf{Paso 2: Multiplicar por el factor integrante 
$\mu = \dfrac{1}{xy(2x+y)}$.}

\[
    \frac{3xy+y^2}{xy(2x+y)} + \frac{x^2+xy}{xy(2x+y)}\,y' = 0
\]

Simplificando cada término:
\[
    M^* = \frac{y(3x+y)}{xy(2x+y)} = \frac{3x+y}{x(2x+y)}
\]
\[
    N^* = \frac{x(x+y)}{xy(2x+y)} = \frac{x+y}{y(2x+y)}
\]

\medskip
\textbf{Paso 3: Verificar exactitud de la ecuación modificada.}

\[
    \frac{\partial M^*}{\partial y} = \frac{-1}{(2x+y)^2} \qquad
    \frac{\partial N^*}{\partial x} = \frac{-1}{(2x+y)^2}
\]
Como son iguales, la ecuación modificada \textbf{es exacta}. $\checkmark$

\medskip
\textbf{Paso 4: Encontrar $F(x,y)$ integrando $N^*$ respecto a $y$.}

\[
    F = \int \frac{x+y}{y(2x+y)}\,dy
\]

Descomponemos en fracciones parciales respecto a $y$:
\[
    \frac{x+y}{y(2x+y)} = \frac{1}{2} \cdot \frac{1}{y} 
    + \frac{1}{2} \cdot \frac{1}{2x+y}
\]

Integrando:
\[
    F = \frac{1}{2}\ln|y| + \frac{1}{2}\ln|2x+y| + g(x)
    = \frac{1}{2}\ln|y(2x+y)| + g(x)
\]

\medskip
\textbf{Paso 5: Determinar $g(x)$ usando $F_x = M^*$.}

\[
    \frac{\partial F}{\partial x} = \frac{1}{2} \cdot \frac{2}{2x+y} 
    + g'(x) = \frac{1}{2x+y} + g'(x)
\]

Igualando a $M^* = \dfrac{3x+y}{x(2x+y)}$:

\[
    \frac{1}{2x+y} + g'(x) = \frac{3x+y}{x(2x+y)}
\]
\[
    g'(x) = \frac{3x+y}{x(2x+y)} - \frac{1}{2x+y} 
    = \frac{3x+y - x}{x(2x+y)} = \frac{2x}{x(2x+y)} \cdot\frac{1}{1}
\]

Simplificando con fracciones parciales:
\[
    g'(x) = \frac{1}{x} \implies g(x) = \ln|x|
\]

\medskip
\textbf{Paso 6: Solución general.}

\[
    F = \ln|x| + \frac{1}{2}\ln|y(2x+y)| = C
\]

Multiplicando por 2 y aplicando propiedades del logaritmo:
\[
    \ln\left|x^2 y(2x+y)\right| = C
\]
\[
    \boxed{x^2 y(2x+y) = C}
\]

\medskip
\textbf{Verificación (computacional).} Se confirmó mediante 
cálculo simbólico que $\dfrac{\partial M^*}{\partial y} = 
\dfrac{\partial N^*}{\partial x} = \dfrac{-1}{(2x+y)^2}$, 
garantizando la exactitud de la ecuación modificada. $\checkmark$

\end{respuesta}

\newpage




